\section{Gauge fixing and fluctuation coordinates}\label{sec:fluctuations}

Throughout this section $V'$ is as in Lemma~\ref{lem:composition}, $W:=V^0$ is the base configuration and $W'=\mathfrak a'(W)=V'$. We use the notation of Section~\ref{sec:adapted}, with the reference $U'$, so that $W_k=W$ and $W_{k+1}=W'$. Configurations on $\T$ are written $V=XW$, with the relative configuration $X$, and $X=e^{\ii Z}$ with the \emph{fluctuation field} $Z$; in the units of step $k$ we have $\delta_k=1$. The integrand of \eqref{eq:rg-map} is a function of $X$ through $V=XW$. By invariance of the Haar measure, $dV(b)=dX(b)$. We write $g=g_k$. The averaged contour variables of a configuration $V$ on $\T$ are those of the step, $V(y,x)$ for $y\in\T'$ and $x\in B'(y)$ (Definition~\ref{def:avg-contour}). Let $\mathfrak S_{k+1}$ be the group of the maps $x\mapsto Rx+v$ with $R\in B_4$ and $v\in\T'$, generated by the translations of $\T'$ and the hyperoctahedral group about the points of $\T'$; it is the group $\mathfrak S'$ of Section~\ref{sec:rg-step}, and $\mathfrak S_{k+1}\subset\mathfrak S_k$.

\subsection{Bonds and blocks}\label{ss:32.1}

\begin{definition}[Crossing bonds and the variables $\mathfrak F$]\label{def:bonds}
We use the following sets of bonds of $\T$.
\begin{itemize}
\item The blocks $B'(y)$, $y\in\T'$, are the cubes of $L^d$ points of $\T$ centred at $y$; they partition $\T$. A bond of $\T$ is \emph{inner} if its endpoints lie in one block, and \emph{crossing} otherwise. A crossing bond $(x,x+e_\mu)$ joins $B'(y)$ to $B'(y+Le_\mu)$.
\item For a bond $c=(y,y+Le_\mu)$ of $\T'$, let $b_0(c)$ be the middle bond of the straight contour $[c]$, from $y+\frac{L-1}2e_\mu$ to $y+\frac{L+1}2e_\mu$. It is the only crossing bond of $[c]$.
\item $\mathfrak F$ is the set of positively oriented bonds of $\T$ other than the bonds $b_0(c)$.
\end{itemize}
\end{definition}

Every $r\in\mathfrak S_{k+1}$ maps blocks to blocks, maps $[c]$ onto $\pm[rc]$ and the middle bond onto the middle bond. So the set of crossing bonds, the set of the $b_0(c)$ and $\mathfrak F$ are invariant under $\mathfrak S_{k+1}$, as sets of unoriented bonds. The contours of a family $G(y,x)$, $x\in B'(y)$, lie in the box spanned by $y$ and $x$, hence in $B'(y)$ (Section~\ref{sec:averaging}). So they consist of inner bonds and contain no bond $b_0(c)$. See Figure~\ref{fig:blocks-sym}.

\begin{figure}[t]
\centering
\begin{tikzpicture}[scale=0.9]
\foreach \bx in {0,3} {
  \fill[fblue!7] (\bx-1.4,-1.4) rectangle (\bx+1.4,1.4);
  \foreach \i in {-1,0,1} \foreach \j in {-1,0} \draw[gray!60] (\bx+\i,\j)--(\bx+\i,\j+1);
  \foreach \i in {-1,0} \foreach \j in {-1,0,1} \draw[gray!60] (\bx+\i,\j)--(\bx+\i+1,\j);
}
\foreach \j in {-1,1} \draw[forange,thick] (1,\j)--(2,\j);
\draw[fred,ultra thick] (1,0)--(2,0);
\draw[fblue,very thick] (0,0)--(1,0)--(1,1);
\draw[fblue,very thick,dashed] (0,0)--(0,1)--(1,1);
\foreach \bx in {0,3} \foreach \i in {-1,0,1} \foreach \j in {-1,0,1} \fill[gray] (\bx+\i,\j) circle (1.6pt);
\fill[black] (0,0) circle (2.6pt);
\fill[black] (3,0) circle (2.6pt);
\node[below left,font=\small] at (0,-0.05) {$y$};
\node[below right,font=\small] at (3,-0.05) {$y+Le_1$};
\node[above right,font=\small] at (1,1) {$x$};
\node[fred,font=\footnotesize] at (1.5,-0.4) {$b_0(c)$};
\node[font=\small,align=left,anchor=west] at (5.2,1.2) {\textcolor{fblue}{blue}: the two contours of $G(y,x)$};
\node[font=\small,align=left,anchor=west] at (5.2,0.4) {\textcolor{forange}{orange}: crossing bonds in $\mathfrak F$};
\node[font=\small,align=left,anchor=west] at (5.2,-0.4) {\textcolor{fred}{red}: $b_0(c)$, eliminated by the constraint};
\node[font=\small,align=left,anchor=west] at (5.2,-1.2) {\textcolor{gray}{thin}: inner bonds, in $\mathfrak F$};
\end{tikzpicture}
\caption{Two centred blocks $B'(y)$, $B'(y+Le_1)$ of $\T$ for $d=2$, $L=3$. The bond $b_0(c)$ is the middle bond of the straight contour $[c]$ from $y$ to $y+Le_1$. The blocks, the crossing bonds and $b_0(c)$ are invariant under the reflections that map the pair of blocks onto itself; the family $G(y,x)$ is mapped onto the family of the image of $x$.}
\label{fig:blocks-sym}
\end{figure}

\subsection{Ba\l aban's exponential gauge fixing}\label{ss:32.2}

The integrand of \eqref{eq:rg-map} is invariant under the group
\[
\mathcal G'=\{u:\T\to G:\ u|_{\T'}=1\}
\]
of Section~\ref{sec:rg-step}, by Lemma~\ref{lem:one-step-basic}(b). Following Ba\l aban \cite[(0.14)--(0.16)]{BalYM4I}, we fix this invariance by exponential weights in the averaged contour variables, not by $\delta$-functions; the latter would define a nonlinear slice. Let
\[
\mathfrak M=\big\{V\ \text{real configuration on }\T:\ |\ell(V(\partial p))|\le\eps_M\ \text{for all plaquettes }p\text{ of }\T\big\},\qquad\eps_M=\frac{r_M}{8L^2},
\]
be a quantitative form of Ba\l aban's set of regular configurations \cite[p.~254]{BalYM4I}; in the units of step $k$, $\mathfrak M=\mathfrak U_k(\eps_M)$. By Lemma~\ref{lem:stokes-avg} and Lemma~\ref{lem:one-step-basic}(a), for $V\in\mathfrak M$ all averaged contour variables $V(y,x)$ are defined and real, their families lie in $\cD_{3L^2\eps_M}\subset\cD_{r_M/2}$, and $\mathfrak a'(V)$ is defined and real. The set $\mathfrak M$ is invariant under real gauge transformations of any size and under $\mathfrak S_{k+1}$ (Section~\ref{sec:averaging}). The base lies in it with a margin: by Lemma~\ref{lem:composition}(a), $W\in\mathfrak U_k(2(C_*+1)L^{-2}\eps'_1)$, and $(C_*+1)\eps'_1\le c_1\le r_M/128$ (Lemma~\ref{lem:one-step-regularity}), so
\begin{equation}\label{eq:base-regular}
|\ell(W(\partial p))|\le2(C_*+1)L^{-2}\eps'_1\le\tfrac18\eps_M\qquad\text{for all plaquettes }p.
\end{equation}

\begin{definition}[Gauge fixing relative to the base]\label{def:gauge-fixing}
For $u\in G^c$ put
\[
\phi(u)=\exp\Big(-\frac1{g^2}\Big[\mathsf N-\tfrac12\operatorname{tr}\big(u+u^{-1}\big)\Big]\Big),\qquad z=z(g)=\int_G\phi(u)\,du .
\]
For a real configuration $W_\circ\in\mathfrak M$ and a configuration $V$ whose averaged contour variables are defined, put
\[
\mathrm{gf}_{W_\circ}(V)=\prod_{y\in\T'}\ \prod_{x\in B'(y)\setminus\{y\}}z^{-1}\,\phi\big(V(y,x)\,W_\circ(y,x)^{-1}\big).
\]
In this part $W_\circ=W$ is the base, and we write $\mathrm{gf}_W$.
\end{definition}

For $u\in G$, $\mathsf N-\frac12\operatorname{tr}(u+u^{-1})=\mathsf N-\operatorname{Re}\operatorname{tr}u=\frac12|u-1|^2$, and at $u=e^{\ii H}$ it equals $f(H):=\mathsf N-\operatorname{tr}\cos H=\frac12\operatorname{tr}H^2+O(|H|^4)$. This is the normalization of the plaquette terms of the Wilson action in Section~\ref{sec:action}, and $\operatorname{tr}H^2=|H|^2$ for $H\in\mathfrak g$. The function $\phi$ is a class function, $0<\phi\le1$ on $G$, and $\phi(u)=1$ only for $u=1$. By Laplace's method, $z(g)=\sigma(0)(2\pi g^2)^{\dim\mathfrak g/2}(1+O(g^2))$ as $g\to0$, with $\sigma$ the density of the Haar measure in exponential coordinates (Section~\ref{subsec:exact-formula}). The strength of the gauge fixing is that of the action, as in \cite[(0.19)]{BalYM4I}. For any fixed $a_{\rm gf}>0$ one may replace $g^{-2}$ in the weight by $a_{\rm gf}g^{-2}$ and compute $z$ at that same strength. The normalized Haar-orbit identity is unchanged; the full function $\Phi_W$, its remainder $G_3$, and the quadratic gauge-fixing term are multiplied by $a_{\rm gf}$. The ensuing coercivity and upper-bound constants may depend on this fixed positive strength, and the normalization changes by a background-independent factor. Ba\l aban's weights contain in addition the restriction $\chi(|U(y,x)-1|<\eps_0)$; here the small-field region is defined by the cutoff of Definition~\ref{def:chi-w} instead, on whose support that factor would be $1$.

For a complex field $Z$ on the bonds of $\T$ and $x\in B'(y)$ put
\[
H_{y,x}(Z)=\ell\big((e^{\ii Z}W)(y,x)\,W(y,x)^{-1}\big),
\]
whenever the averaged contour variable is defined and $|(e^{\ii Z}W)(y,x)\,W(y,x)^{-1}-1|<1$, the domain of $\ell$ (Section~\ref{ss:21.2}). The proof of Lemma~\ref{lem:exp-gauge}(c) verifies both conditions on its small complex $Z$-ball, by Lemma~\ref{lem:group-average}(a),(e). For complex fields we use the bilinear pairing $\ip AB=\sum_{y\in\T'}\sum_{x\in B'(y)}\operatorname{tr}(A(y,x)B(y,x))$ on functions of the pairs $(y,x)$; for real fields $\ip AA=\|A\|^2$. Recall $\mathfrak G_WZ(y,x)=\lambda_W(y,x)[Z]$ from Definition~\ref{def:avg-transport}, in the units of step $k$.

\begin{lemma}[Exponential gauge fixing]\label{lem:exp-gauge}
Let $W_\circ\in\mathfrak M$ be real.
\begin{enumerate}[label=(\alph*)]
\item \emph{(Faddeev--Popov identity.)} Let $f$ be an integrable function of real configurations on $\T$ with $f(V^u)=f(V)$ for all $u\in\mathcal G'$, and $f=0$ outside $\mathfrak M$. Then
\[
\int\prod_bdV(b)\ f(V)=\int\prod_bdV(b)\ \mathrm{gf}_{W_\circ}(V)\,f(V).
\]
\item \emph{(Invariance.)} For $V\in\mathfrak M$ and every real gauge transformation $u$ on $\T$, of any size, $\mathrm{gf}_{W_\circ^u}(V^u)=\mathrm{gf}_{W_\circ}(V)$. For $r\in\mathfrak S_{k+1}$, $\mathrm{gf}_{rW_\circ}(rV)=\mathrm{gf}_{W_\circ}(V)$.
\item \emph{(Expansion.)} For $V\in\mathfrak M$, $0<\mathrm{gf}_{W_\circ}(V)\le z^{-|\T\setminus\T'|}$, with equality if and only if $V(y,x)=W_\circ(y,x)$ for all $y\in\T'$, $x\in B'(y)$. There is $c_0>0$, depending only on $G$, such that for complex $Z$ with $|Z|_\infty\le c_0L^{-1}$ the averaged contour variables of $e^{\ii Z}W_\circ$ are defined, and
\[
\mathrm{gf}_{W_\circ}(e^{\ii Z}W_\circ)=z^{-|\T\setminus\T'|}e^{-g^{-2}\Phi_{W_\circ}(Z)},
\]
\[
\Phi_{W_\circ}(Z)=\sum_{y\in\T'}\sum_{x\in B'(y)}f\big(H_{y,x}(Z)\big)=\tfrac12\ip{\mathfrak G_{W_\circ}Z}{\mathfrak G_{W_\circ}Z}+G_3(Z),
\]
with $H_{y,x}$ computed at $W_\circ$. Here $G_3=\sum_{y\in\T'}G_{3,y}$, where $G_{3,y}$ depends only on $Z$ on the inner bonds of $B'(y)$, is holomorphic, real for real $Z$, vanishes to third order at $Z=0$, and satisfies $|G_{3,y}(Z)|\le C_L\max_b|Z(b)|^3$ for $|Z|_\infty\le\frac12c_0L^{-1}$, the maximum over the bonds of $B'(y)$. In particular $\Phi_{W_\circ}$ depends only on $Z$ on inner bonds, and $|G_3(Z)|\le C_L|Z|_\infty\|Z\|^2$.
\end{enumerate}
\end{lemma}

\begin{proof}\hasdetail{lem:exp-gauge}
(a) For real $u$ of any size and $V\in\mathfrak M$, \eqref{eq:avg-contour-gauge} gives $V^u(y,x)=u(y)V(y,x)u(x)^{-1}$. For $u\in\mathcal G'$ this is $V(y,x)u(x)^{-1}$, since $u(y)=1$. Fix $V\in\mathfrak M$. Each $x\in\T\setminus\T'$ lies in exactly one block $B'(y)$, and is not its centre. The map $u\mapsto V(y,x)u^{-1}W_\circ(y,x)^{-1}$ preserves the normalized Haar measure. Hence, with $du=\prod_{x\in\T\setminus\T'}du(x)$,
\[
\int_{\mathcal G'}\mathrm{gf}_{W_\circ}(V^u)\,du=\prod_{y\in\T'}\prod_{x\in B'(y)\setminus\{y\}}z^{-1}\int_G\phi\big(V(y,x)u(x)^{-1}W_\circ(y,x)^{-1}\big)\,du(x)=1 .
\]
Put $\mathrm{gf}_{W_\circ}=0$ outside $\mathfrak M$. Since $\mathfrak M$ is $\mathcal G'$-invariant, $\mathrm{gf}_{W_\circ}(V^u)f(V^u)=0$ for $V\notin\mathfrak M$. The map $V\mapsto V^u$ acts on each bond by a left and a right translation, so it preserves $\prod_bdV(b)$. Since $0\le\mathrm{gf}_{W_\circ}\le z^{-|\T\setminus\T'|}$ and $f$ is integrable, Fubini's theorem applies:
\[
\int dV\,\mathrm{gf}_{W_\circ}f=\int_{\mathcal G'}du\int dV\,\mathrm{gf}_{W_\circ}(V^u)f(V^u)=\int dV\,f(V)\int_{\mathcal G'}\mathrm{gf}_{W_\circ}(V^u)\,du=\int dV\,f(V).
\]
No Faddeev--Popov determinant appears, since the $\mathcal G'$-orbits are integrated with the Haar measure.

(b) By \eqref{eq:avg-contour-gauge}, $V^u(y,x)W_\circ^u(y,x)^{-1}=u(y)\big(V(y,x)W_\circ(y,x)^{-1}\big)u(y)^{-1}$, and $\phi$ is a class function. For $r\in\mathfrak S_{k+1}$, $(rV)(y,x)=V(r^{-1}y,r^{-1}x)$ (Section~\ref{sec:averaging}), and $r$ maps the set of pairs $(y,x)$ with $y\in\T'$, $x\in B'(y)\setminus\{y\}$ bijectively onto itself, since it maps $\T'$ onto $\T'$ and $B'(y)$ onto $B'(ry)$.

(c) The first statement follows from $\phi=e^{-|u-1|^2/2g^2}$ on $G$. Let $\mathbf U=\{W_\circ(\Gamma)\}_{\Gamma\in G(y,x)}$, which lies in $\cD_{r_M/2}$. For $\Gamma\in G(y,x)$ write $(e^{\ii Z}W_\circ)(\Gamma)=e^{\ii\xi_\Gamma}W_\circ(\Gamma)$, where $e^{\ii\xi_\Gamma}=(e^{\ii Z})_{W_\circ}(\Gamma)$ is a product of at most $2(L-1)$ factors $e^{\ii R(W_\circ(\cdot))Z(b)}$ with unitary $R(W_\circ(\cdot))$. By Lemma~\ref{lem:log-estimates}(c), $\xi_\Gamma$ is holomorphic in $Z$ with $|\xi_\Gamma|\le4L|Z|_\infty\le r_M/16$, if $c_0$ is small. By Lemma~\ref{lem:group-average}(a),(e) (and its proof), the family $\{e^{\ii\xi_\Gamma}W_\circ(\Gamma)\}$ lies in $\cD_{r_M}$, the averaged contour variable $(e^{\ii Z}W_\circ)(y,x)=\Av(\{e^{\ii\xi_\Gamma}W_\circ(\Gamma)\})$ is defined and holomorphic in $Z$, and
\[
|H_{y,x}(Z)|=\big|\ell\big(\Av(\{e^{\ii\xi_\Gamma}W_\circ(\Gamma)\})\Av(\mathbf U)^{-1}\big)\big|\le C_M\max_\Gamma|\xi_\Gamma|\le4C_ML|Z|_\infty .
\]
For $c_0$ small this is at most $1$. Since $e^{\ii\ell(h)}=h$, we have $V(y,x)W_\circ(y,x)^{-1}=e^{\ii H_{y,x}(Z)}$ for $V=e^{\ii Z}W_\circ$, and $\phi(e^{\ii H})=e^{-g^{-2}f(H)}$. This gives the formula with $\Phi_{W_\circ}=\sum f(H_{y,x})$. The function $H_{y,x}$ depends only on $Z$ on the bonds of the box spanned by $y$ and $x$, which are inner bonds of $B'(y)$. By Lemma~\ref{lem:avg-transport}(a), $H_{y,x}(sZ)=s\,\mathfrak G_{W_\circ}Z(y,x)+O(s^2)$. Since $f(H)=\frac12\operatorname{tr}H^2+O(|H|^4)$, the second-order part of $\Phi_{W_\circ}$ is $\frac12\ip{\mathfrak G_{W_\circ}Z}{\mathfrak G_{W_\circ}Z}$, and
\[
G_{3,y}(Z)=\sum_{x\in B'(y)}\Big[f\big(H_{y,x}(Z)\big)-\tfrac12\operatorname{tr}\big((\mathfrak G_{W_\circ}Z(y,x))^2\big)\Big]
\]
vanishes to third order. Reality follows from that of $H_{y,x}$ for real $Z$. For the bound let $m$ be the maximum of $|Z(b)|$ over the bonds of $B'(y)$, and $t\mapsto G_{3,y}(tZ)$ for $|t|\le R:=c_0/(Lm)$. There $|H_{y,x}|\le4C_Mc_0\le1$ and $|\mathfrak G_{W_\circ}(tZ)(y,x)|\le C$ (Lemma~\ref{lem:avg-transport}(c)), so $|G_{3,y}(tZ)|\le CL^d$. The maximum principle for $G_{3,y}(tZ)/t^3$ gives $|G_{3,y}(Z)|\le CL^dR^{-3}=CL^{d+3}c_0^{-3}m^3$. Finally $m^3\le|Z|_\infty\sum_{b}|Z(b)|^2$, the sum over the bonds of $B'(y)$, and every bond lies in at most one block.
\end{proof}

Part~(a) is \cite[(0.15)--(0.16)]{BalYM4I}. Ba\l aban's weights are relative to $1$, i.e.\ $W_\circ(y,x)=1$. The real transformation $u\in\mathcal G'$ with $u(x)=W(y,x)$ for $x\in B'(y)$ satisfies $W^u(y,x)=1$, so by~(b) $\mathrm{gf}_W(V)=\mathrm{gf}_{W^u}(V^u)$ is Ba\l aban's weight evaluated at $V^u$. This is the gauge in which Ba\l aban expands the fluctuation integral, with the expansion~(c) of the gauge-fixing term \cite[(2.3)--(2.5)]{BalYM4I}. Relative to the base the gauge fixing needs no gauge condition on $W$, so the local gauges of Parts~VIII--X remain available.

\begin{definition}[Fluctuation cutoff]\label{def:chi-w}
For $w\le\frac12c_G$ and a configuration $V$ on $\T$ put $\chi^{\rm w}(V)=1$ if for every bond $b$ of $\T$ there is $Z(b)\in\mathfrak g$ with $|Z(b)|\le w$ and $V(b)=e^{\ii Z(b)}W(b)$, and $\chi^{\rm w}(V)=0$ otherwise.
\end{definition}

For $w\le\frac12c_G$ the field $Z(b)=\ell(V(b)W(b)^{-1})$ is unique (Lemma~\ref{lem:log-estimates}(a),(d)). The cutoff is a product over all bonds, including the $b_0(c)$; this is Ba\l aban's characteristic function \cite[(2.9)]{BalYM4I}. It depends on $V'$ through $W$. It is not $\mathcal G'$-invariant, but it is covariant: $|\ell(V(b)W(b)^{-1})|$ does not change if $(V,W)$ is replaced by $(V^u,W^u)$ for a real gauge transformation $u$, since then $V(b)W(b)^{-1}$ is conjugated by $u(b_-)$, or by $(rV,rW)$ for $r\in\mathfrak S_{k+1}$, since on reversed bonds $V(-b)W(-b)^{-1}=W(b)^{-1}\big(V(b)W(b)^{-1}\big)^{-1}W(b)$. Let $w\le\frac18\eps_M$. On the support of $\chi^{\rm w}$, write $V(\partial p)$ as a product of four conjugates of the $e^{\pm\ii Z(b)}$ and of $W(\partial p)$, by moving the factors $W(b)$ to the right. By \eqref{eq:base-regular} and Lemma~\ref{lem:log-estimates}(b), every plaquette satisfies $|\ell(V(\partial p))|\le|\ell(W(\partial p))|+4w$. Hence
\begin{equation}\label{eq:chi-w-support}
\operatorname{supp}\chi^{\rm w}\subset\mathfrak U_k\big(2(C_*+1)L^{-2}\eps'_1+4w\big)\subset\mathfrak M\qquad\big(w\le\tfrac18\eps_M\big).
\end{equation}

\subsection{The gauge-fixed step}\label{ss:32.3}

\begin{definition}[Gauge-fixed step]\label{def:gf-step}
For a function $f$ of real configurations on $\T$ that vanishes outside $\mathfrak M$ put
\[
T^{\rm gf}(f)(V')=\int\prod_bdV(b)\ \prod_c\delta\big(\mathfrak a'(V)(c)\,V'(c)^{-1}\big)\ \mathrm{gf}_W(V)\,f(V),
\]
with $W$ the base of $V'$, and the measures and the $\delta$-function of \eqref{eq:rg-map}.
\end{definition}

The integrand may depend on $V'$ through the base $W$ and the reference $U'$; we then write $f=f_{W,U'}$. We call a function $\chi=\chi_{W,U'}$ of configurations $V$ on $\T$, with $0\le\chi\le1$, a \emph{covariant cutoff} if
\begin{enumerate}[label=(C\arabic*)]
\item $\chi_{W^{u|_\T},U'^u}(V^{u|_\T})=\chi_{W,U'}(V)$ for every real gauge transformation $u$ on $\Tf$;
\item $\chi_{rW,rU'}(rV)=\chi_{W,U'}(V)$ for every $r\in\mathfrak S_{k+1}$.
\end{enumerate}
By the paragraph after Definition~\ref{def:chi-w}, $\chi^{\rm w}$ is a covariant cutoff. The identities below between integrals containing $\delta$-functions are first proved with $\delta$ replaced by a continuous class function that is invariant under inversion and concentrated near $1$, for instance a function of $|\ell(u)|$ supported near $1$, and then hold in the limit, as identities of measures. Where $\mathfrak a'$ is a submersion they hold also pointwise for the fibre densities of Section~\ref{sec:rg-step}, at every $V'$ at which the fibre integrals converge absolutely: the substitutions used below are isometries of the bond variables that map fibres of $\mathfrak a'$ onto fibres, and the fibre measure is invariant under $\mathcal G'$. Both invariances are needed: the proofs use $\delta(a^{-1}b)=\delta(ab^{-1})$, and a class function need not be invariant under $u\mapsto u^{-1}$ (for instance on $SU(3)$).

\begin{lemma}[Invariance of the gauge-fixed step]\label{lem:gf-step-invariance}
Let $\rho$ be a gauge-invariant function of real configurations on $\T$, let $\chi$ be a covariant cutoff, and suppose that $\chi\rho$ vanishes outside $\mathfrak M$.
\begin{enumerate}[label=(\roman*)]
\item $T^{\rm gf}(\chi\rho)(V')$ does not depend on the choice of $U'$ in the minimal orbit of $V'$.
\item $T^{\rm gf}(\chi\rho)(V'^{u'})=T^{\rm gf}(\chi\rho)(V')$ for every real gauge transformation $u'$ on $\T'$.
\item For every $r\in\mathfrak S_{k+1}$, $T^{\rm gf}(\chi\cdot r\rho)=r\,T^{\rm gf}(\chi\rho)$, with the actions $(r\rho)(V)=\rho(r^{-1}V)$ and $(rF)(V')=F(r^{-1}V')$ of Section~\ref{sec:rg-step}; that is, $T^{\rm gf}(\chi\cdot r\rho)(rV')=T^{\rm gf}(\chi\rho)(V')$, where on each side $\chi$ is taken at the base and the reference of the argument. In particular, if $\rho(rV)=\rho(V)$ for all $V$, then $T^{\rm gf}(\chi\rho)(rV')=T^{\rm gf}(\chi\rho)(V')$.
\item On the set of the $V'$ as in Lemma~\ref{lem:composition}, $T^{\rm gf}(f)=T(f)$ for every integrable $f$ with $f(V^u)=f(V)$ for $u\in\mathcal G'$ that vanishes outside $\mathfrak M$. Consequently, let $w\le\frac18\eps_M$ and $2(C_*+1)L^{-2}\eps'_1+4w\le\tilde h$, so that $\operatorname{supp}\chi^{\rm w}\subset\mathfrak M\cap\mathfrak U_k(\tilde h)$. Then, on the same set, for every covariant cutoff $\chi$ and every gauge-invariant $\rho\ge0$ on $\mathfrak U_k(\tilde h)$,
\[
T\big(\one_{\mathfrak M\cap\mathfrak U_k(\tilde h)}\rho\big)-T^{\rm gf}(\chi^{\rm w}\chi\rho)=T^{\rm gf}\big((1-\chi^{\rm w}\chi)\one_{\mathfrak M\cap\mathfrak U_k(\tilde h)}\rho\big)\ge0 .
\]
\end{enumerate}
\end{lemma}

\begin{proof}\hasdetail{lem:gf-step-invariance}
The proofs of (i)--(iii) substitute $V=\tilde V^{v}$ or $V=r\tilde V$ in the integral. These substitutions preserve $\prod_bdV(b)$: a gauge transformation acts on each bond by a left and a right translation, and $r$ permutes the bonds, possibly reversing them, which preserves the Haar measure. They also preserve $\mathfrak M$.

(i) Another element of the minimal orbit is $U'^v$ with $v$ real on $\Tf$ and $v|_{\T'}=1$. By Lemma~\ref{lem:one-step-basic}(b), its base is $W^{\tilde v}$ with $\tilde v=v|_\T\in\mathcal G'$. Substitute $V=\tilde V^{\tilde v}$. Then $\mathfrak a'(V)=\mathfrak a'(\tilde V)$, since $\tilde v|_{\T'}=1$. Moreover $\mathrm{gf}_{W^{\tilde v}}(\tilde V^{\tilde v})=\mathrm{gf}_W(\tilde V)$ by Lemma~\ref{lem:exp-gauge}(b), $\chi_{W^{\tilde v},U'^v}(\tilde V^{\tilde v})=\chi_{W,U'}(\tilde V)$ by (C1), and $\rho(\tilde V^{\tilde v})=\rho(\tilde V)$. So the integral with $U'^v$ equals the integral with $U'$.

(ii) Let $u$ be a real gauge transformation on $\Tf$ with $u|_{\T'}=u'$, and $\tilde u=u|_\T$. The configuration $V'^{u'}$ satisfies the hypotheses of Lemma~\ref{lem:composition}, since its curvature is that of $V'$. By Theorem~\ref{thm:variational}(g) and Theorem~\ref{thm:balaban}(a) for the $(k+1)$-step problem, $U'^u$ lies in the minimal orbit of $V'^{u'}$, and by Lemma~\ref{lem:one-step-basic}(b) its base is $W^{\tilde u}$. By (i) we may use this representative. Substitute $V=\tilde V^{\tilde u}$. By Lemma~\ref{lem:one-step-basic}(b),
\[
\mathfrak a'(\tilde V^{\tilde u})(c)\,V'^{u'}(c)^{-1}=u'(c_-)\,\mathfrak a'(\tilde V)(c)V'(c)^{-1}\,u'(c_-)^{-1},
\]
and $\delta$ is a class function. The other factors are invariant as in (i), now with $\tilde u$ in place of $\tilde v$.

(iii) By Proposition~\ref{prop:symmetry}(b) and its proof, $rV'$ satisfies the hypotheses of Lemma~\ref{lem:composition}, and $rU'$ lies in the minimal orbit of $rV'$. Since $\mathfrak S_{k+1}\subset\mathfrak S_k$, Lemma~\ref{lem:one-step-basic} gives $(rU')^{(k)}=rW$. By (i) we may use the reference $rU'$ and the base $rW$ at $rV'$. The hypotheses of the lemma hold with $r\rho$ in place of $\rho$: $r\rho$ is gauge invariant (Proposition~\ref{prop:symmetry}(a)), and $\chi_{rW,rU'}\cdot r\rho=r\big(\chi_{W,U'}\rho\big)$ by (C2), which vanishes outside $r\mathfrak M=\mathfrak M$. Substitute $V=r\tilde V$ in the integral for $T^{\rm gf}(\chi\cdot r\rho)(rV')$. By Lemma~\ref{lem:one-step-basic}(d), $\mathfrak a'(r\tilde V)(c)=\mathfrak a'(\tilde V)(r^{-1}c)$, and $(rV')(c)=V'(r^{-1}c)$. As $c$ runs over the positively oriented bonds of $\T'$, $r^{-1}c$ runs over all bonds of $\T'$, each once up to orientation. If $r^{-1}c=-c''$, then $\mathfrak a'(\tilde V)(-c'')V'(-c'')^{-1}=\mathfrak a'(\tilde V)(c'')^{-1}V'(c'')$ by Lemma~\ref{lem:reversal}, and $\delta(a^{-1}b)=\delta(ab^{-1})$, since $\delta$ is invariant under inversion and conjugation. So the product of the $\delta$-functions is unchanged. Moreover $\mathrm{gf}_{rW}(r\tilde V)=\mathrm{gf}_W(\tilde V)$ by Lemma~\ref{lem:exp-gauge}(b), $\chi_{rW,rU'}(r\tilde V)=\chi_{W,U'}(\tilde V)$ by (C2), and $(r\rho)(r\tilde V)=\rho(\tilde V)$. So the integral equals $T^{\rm gf}(\chi\rho)(V')$. If $\rho\circ r=\rho$, then $r\rho=\rho$, which gives the last claim.

(iv) Let $\varphi$ be a bounded Borel function of configurations on $\T'$ that vanishes outside the set of the $V'$ as in Lemma~\ref{lem:composition}. The function $V\mapsto\varphi(\mathfrak a'(V))f(V)$ is integrable, $\mathcal G'$-invariant by Lemma~\ref{lem:one-step-basic}(b), and vanishes outside $\mathfrak M$. The base $W$, taken at $V'=\mathfrak a'(V)$, is constant on $\mathcal G'$-orbits, and in the proof of Lemma~\ref{lem:exp-gauge}(a) the orbit integral of $\mathrm{gf}_{W_\circ}$ equals $1$ for every $W_\circ$. So that proof, with $W_\circ$ the base on each orbit, gives $\int\prod_bdV(b)\,\varphi(\mathfrak a'(V))\,\mathrm{gf}_W(V)f(V)=\int\prod_bdV(b)\,\varphi(\mathfrak a'(V))f(V)$, that is, $T^{\rm gf}(f)=T(f)$ as measures on that set. If $\mathfrak a'$ is a submersion near the support of $f$ and the fibre integral of $|f|$ at $V'$ is finite, the same argument, run on the fibre, whose measure is $\mathcal G'$-invariant, gives the identity of the fibre densities at $V'$. The inclusion $\operatorname{supp}\chi^{\rm w}\subset\mathfrak M\cap\mathfrak U_k(\tilde h)$ is \eqref{eq:chi-w-support}. The function $\one_{\mathfrak M\cap\mathfrak U_k(\tilde h)}\rho$ is $\mathcal G'$-invariant, since $\mathfrak M$ and $\mathfrak U_k(\tilde h)$ are invariant under real gauge transformations. So $T(\one_{\mathfrak M\cap\mathfrak U_k(\tilde h)}\rho)=T^{\rm gf}(\one_{\mathfrak M\cap\mathfrak U_k(\tilde h)}\rho)$. By the inclusion, $\chi^{\rm w}\chi\rho=\chi^{\rm w}\chi\one_{\mathfrak M\cap\mathfrak U_k(\tilde h)}\rho$. Subtracting gives the identity. The right side is nonnegative, since $\delta\ge0$, $\mathrm{gf}_W\ge0$, $0\le\chi^{\rm w}\chi\le1$ and $\rho\ge0$.
\end{proof}

Lemma~\ref{lem:gf-step-invariance} describes the \emph{modified model}. The small-field step is $T^{\rm gf}(\chi^{\rm w}\chi_\xi\rho)$, with a white-noise cutoff $\chi_\xi$ (Section~\ref{sec:cluster}). The cutoffs act on gauge-fixed variables, so this is not of the form $T(\chi\rho)$ with a $\mathcal G'$-invariant $\chi$. Parts~(i)--(iii) show that it is nevertheless a renormalization group map on gauge-invariant densities that commutes with the lattice symmetries: $T^{\rm gf}(\chi^{\rm w}\chi_\xi\cdot r\rho)=r\,T^{\rm gf}(\chi^{\rm w}\chi_\xi\rho)$ for $r\in\mathfrak S_{k+1}$, by~(iii), since $\chi^{\rm w}\chi_\xi$ is a covariant cutoff (see below). Part~(iv) shows that the discarded part, the difference between the step without cutoffs and the small-field step, is nonnegative. In Section~\ref{sec:cluster} the white-noise cutoff is taken bondwise on $\mathfrak F$ in the coordinates of this section:
\[
\chi_\xi(V)=\prod_{b\in\mathfrak F}\one\big(|\xi_V(b)|\le p_0\big),\qquad\xi_V=g^{-1}C_k^{-1/2}Z_V,\quad Z_V(b)=\ell\big(V(b)W(b)^{-1}\big),\ b\in\mathfrak F,
\]
for $V$ in the support of $\chi^{\rm w}$, with $C_k=C_k(W,U')$ the covariance of Corollary~\ref{cor:covariance}. By Corollary~\ref{cor:covariance}(a),(b), $\xi_{V^{\tilde u}}=R(\tilde u)\xi_V$ and $\xi_{rV}=r_W\xi_V$, with the data transformed accordingly, so $\chi^{\rm w}\chi_\xi$ is a covariant cutoff.

\subsection{Elimination of the constraint}\label{ss:32.4}

For a complex field $Z$ on the bonds of $\T$ and a bond $c$ of $\T'$ put
\[
\ell_c(Z)=\ell\big(\overline X(c)\big),\qquad\overline X=\mathfrak a'(e^{\ii Z}W)\,W'^{-1},
\]
and let $l^{\rm pl}$ be its linearization, $(l^{\rm pl}Z)(c)=\frac{d}{ds}\big|_{s=0}\ell_c(sZ)$; ``pl'' refers to the plain average \eqref{eq:one-step}. The constraint in \eqref{eq:rg-map} is $\ell_c(Z)=0$ for all $c$. Let $\mathsf a:=2(C_*+1)\eps'_1$. By Lemma~\ref{lem:composition} and Proposition~\ref{prop:sharp-curvature}, it bounds the curvature both of $W$ and of the fine reference $U'$, in the units of step $k+1$. The former is what Lemma~\ref{lem:one-step-expansion} needs: by the paragraph after its proof, it applies with $a=\mathsf a$ and $W_{j-1}=W$, which is real and lies in $\mathfrak U_k(\mathsf a)\subset\mathfrak U_k(2\mathsf a)$. In the units of step $k+1$ we write $p'=p_{k+1}=l_{k+1}-\delta_{k+1}D_{W'}\gamma_{k+1}$ for the operator of Lemma~\ref{lem:one-step-expansion}(f) at the reference $U'$, and $q'=q_{k+1}$, $r'=r_{k+1}$, $\gamma'_{\rm lin}=\gamma_{k+1}$.

\begin{lemma}[Local elimination]\label{lem:eliminate}
There are $C,c>0$, depending only on $d$ and $G$, such that the following holds with $\rho_L=cL^{-d-1}$ and $r_L=c^2L^{-2d-1}$, provided $L^{d}\mathsf a\le c$. Moreover $r_L\le\frac18\eps_M$.
\begin{enumerate}[label=(\alph*)]
\item $\ell_c$ is analytic for complex $Z$ with $|Z|\le cL^{-1}$ on the bonds with both endpoints in $B'(c_-)\cup B'(c_+)$, and depends only on these values. Among the bonds $b_0(c')$ it depends only on $Z(b_0(c))$. Moreover $l^{\rm pl}=Lp'$, and $\ell_c(Z)=(l^{\rm pl}Z)(c)+O(L^2|Z|^2)$.
\item $\partial\ell_c/\partial Z(b_0(c))=L^{1-d}R\big(W([c_-,b_0(c)_-])\big)\big(1+\mathfrak r_c(Z)\big)$ with $|\mathfrak r_c(Z)|\le C(L^d\mathsf a+L^{d+1}|Z|)$.
\item For complex $Z$ on $\mathfrak F$ with $|Z|\le r_L$ there is exactly one $\zeta_c(Z)$ with $|\zeta_c(Z)|\le\rho_L$ such that $\ell_c=0$ when $Z(b_0(c))=\zeta_c(Z)$. It is analytic, real for real $Z$, depends only on $Z$ on the bonds of (a), and satisfies $|\zeta_c(Z)|\le CL^d\max|Z|$, the maximum running over those bonds.
\item Let $\hat Z(Z)$ be the field equal to $Z$ on $\mathfrak F$ and to $\zeta_c(Z)$ on $b_0(c)$, and write $\hat Z=\hat Z^1Z+\hat Z^2(Z)+\hat Z^{\ge3}(Z)$ (Taylor expansion). Then $\hat Z^1$ is a linear bijection from fields on $\mathfrak F$ onto
\[
\mathcal Z=\ker l^{\rm pl}=\{B:\ l^{\rm pl}B=0\},
\]
and $\|Z\|\le\|\hat Z^1Z\|\le CL^{d}\|Z\|$. The maps $\hat Z^2$, $\hat Z^{\ge3}$ are supported on the bonds $b_0(c)$ and satisfy $|\hat Z^2(Z)(b_0(c))|\le CL^{3d+1}\max|Z|^2$, the maximum over the bonds of (a).
\end{enumerate}
\end{lemma}

\begin{proof}\hasdetail{lem:eliminate}
(a), (b) We apply Lemma~\ref{lem:one-step-expansion}(f) in the units of step $k+1$, with $a=\mathsf a\le c_2$; its constants $c_2$ and $C$ depend only on $d$ and $G$. There $\delta_k=L^{-1}$, $\delta_{k+1}=1$, and $e^{\ii Z}=e^{\ii\delta_kB}$ with $B=LZ$, so
\[
\ell_c(Z)=(p'B)(c)+c^p_{k+1}(B)(c)=L(p'Z)(c)+c^p_{k+1}(LZ)(c),
\]
for $|B|\le c_2$, i.e.\ $|Z|\le c_2L^{-1}$, with the bounds of Lemma~\ref{lem:one-step-expansion}(c) for $c^p_{k+1}$. This gives $l^{\rm pl}=Lp'$, the expansion, the analyticity and the locality. A bond $b_0(c')$ has both endpoints in $B'(c_-)\cup B'(c_+)$ only if $c'=\pm c$.

Next consider the coefficient of $Z(b_0(c))$ in $L(p'Z)(c)$. By Lemma~\ref{lem:one-step-expansion}(f), $\gamma'_{\rm lin}Z(c_\pm)$ does not depend on $Z(b_0(c))$, and the coefficient of $Z(b_0(c))$ in $(q'Z)(c)$ is $L^{-d}R(W([c_-,b_0(c)_-]))$. The reason is that $[x,x(c)]$ contains $b_0(c)$ only for the $L$ points $x$ on the axis of $c$, for which $G(c_-,x)$ consists of the straight contour $[c_-,x]$. So the coefficient in $L(q'Z)(c)$ is $L^{1-d}R(W([c_-,b_0(c)_-]))$. The bound on $r'$ in Lemma~\ref{lem:one-step-expansion}(b) and the Lipschitz bound for $c^p_{k+1}$ in~(c), applied to $B=LZ$, change it by at most $CL\mathsf a+CL^2|Z|$, i.e.\ relatively by $C(L^d\mathsf a+L^{d+1}|Z|)$. The factor $L^{d+1}$ is genuine. The coefficient $L^{1-d}$ is what remains after a cancellation of terms of order one: the $L$ averaged loops through points of the axis are identically $1$, each of the other $L^d-L$ loops contributes $-L^{-d}$ through $Y([c])^{-1}$, and the outer factor $Y([c])$ contributes $+1$. The non-abelian corrections act on these terms of order one.

(c) Put $M_c=L^{1-d}R(W([c_-,b_0(c)_-]))$ and consider $\Phi(z)=z-M_c^{-1}\ell_c(Z,z)$, where $z$ is the value on $b_0(c)$. By (b), $|\Phi'(z)|\le\frac12$ on $|z|\le\rho_L$ if $c$ is small, since there $L^{d+1}|z|\le c$. By (a), $|\ell_c(Z,0)|\le CL\max|Z|$, since $|p'Z|\le C\max|Z|$ by Lemma~\ref{lem:one-step-expansion}(b). So $|\Phi(0)|\le C_0L^d\max|Z|$, with $C_0$ depending only on $d$ and $G$. For $|Z|\le r_L=c^2L^{-2d-1}$ this is at most $C_0c\,\rho_L\le\frac12\rho_L$ for $c\le1/(2C_0)$. (With the same constant $c$ in $\rho_L$ and $r_L$ the contraction would not close.) So $\Phi$ is a contraction of $\{|z|\le\rho_L\}$ into itself, and its fixed point is the unique zero there. Analyticity follows as in Lemma~\ref{lem:linearizing}. Reality follows from uniqueness, since $\overline X\in G$ for real $Z$. Finally, for $c\le r_M/64$ we have $r_L\le cL^{-2}\le r_M/(64L^2)=\frac18\eps_M$.

(d) Differentiating $\ell_c(\hat Z(Z))=0$ at $Z=0$ gives $(l^{\rm pl}\hat Z^1Z)(c)=0$ for all $c$, so $\hat Z^1Z\in\mathcal Z$. Conversely, $B\in\mathcal Z$ is determined by its values on $\mathfrak F$: by (a), $(l^{\rm pl}B)(c)$ involves among the bonds $b_0(\cdot)$ only $b_0(c)$, with the invertible coefficient of (b) at $Z=0$, so the equation $(l^{\rm pl}B)(c)=0$ fixes $B(b_0(c))$. This gives the bijection. The bound on $\hat Z^1$ follows from $|\zeta^{1}_c(Z)|^2\le CL^{2d}\max|Z|^2\le CL^{2d}\sum|Z(b)|^2$, the sum over the bonds of (a), since each bond lies in a bounded number of these sets. The bound on $\hat Z^2$ follows from Cauchy's estimate, since $|\zeta_c|\le\rho_L$ on $|Z|\le r_L$.
\end{proof}

At the trivial base $W=1$ (and $U'=1$) the averaged contour variables linearize to the means of the line integrals (Lemma~\ref{lem:avg-transport}(b)), and $l^{\rm pl}$ becomes
\begin{equation}\label{eq:q-plain}
(q^{\rm pl}B)(c)=L^{-d}\sum_{x\in B'(c_-)}\Big(\mathfrak G_1B(c_-,x)+B([x,x(c)])-\mathfrak G_1B(c_+,x(c))\Big),
\end{equation}
where $B([x,x(c)])$ is the sum of $B$ over the $L$ bonds of the straight segment. This follows by linearizing \eqref{eq:one-step} at $Y=1$ with Lemma~\ref{lem:group-average}(c), since $L^{-d}\sum_xB([c])=B([c])$ cancels against the outer factor. Equivalently $q^{\rm pl}=L(q-D'\tilde\gamma)$ in the units of step $k+1$, with $q$ the centred tube average, $\tilde\gamma B(y)=L^{-d-1}\sum_{x\in B'(y)}\mathfrak G_1B(y,x)$ and $D'\lambda(c)=\lambda(c_+)-\lambda(c_-)$ (Lemma~\ref{lem:one-step-expansion}(f) and the remark after it). The coefficient of $B(b_0(c))$ in $(q^{\rm pl}B)(c)$ is $L^{1-d}$.

\subsection{The exact formula}\label{subsec:exact-formula}

Let $\sigma$ be the density of the normalized Haar measure in exponential coordinates: $d(e^{\ii Z})=\sigma(Z)\,dZ$ for $|Z|<c_G$, where $dZ$ is the Lebesgue measure on $\mathfrak g$ given by the trace form. It is analytic, positive and $\operatorname{Ad}$-invariant near $0$, and $\sigma(Z)/\sigma(0)=\det\big(\sinh(\tfrac{\ii}{2}\operatorname{ad}Z)/(\tfrac{\ii}{2}\operatorname{ad}Z)\big)=1+O(|Z|^2)$. The $\delta$-function of $G$ is then $\delta(e^{\ii w})=\delta(w)/\sigma(w)$ near $1$.

\begin{proposition}[One step in fluctuation coordinates]\label{prop:exact-step}
Let $V'$ be as in Lemma~\ref{lem:composition}, with $L^d\mathsf a\le c$, and let $w\le r_L$ (so $w\le\rho_L$ and $w\le\frac18\eps_M$). Let $\rho$ be a function of real configurations on $\T$, measurable and bounded on the support of $\chi^{\rm w}$. Then $\chi^{\rm w}\rho$ vanishes outside $\mathfrak M$, which lies in the interior of $\operatorname{dom}\mathfrak a'$, $\mathfrak a'$ is a submersion at every point where $\chi^{\rm w}=1$, and the density of $T^{\rm gf}(\chi^{\rm w}\rho)$ at $V'$, defined in Section~\ref{sec:rg-step}, is
\[
T^{\rm gf}(\chi^{\rm w}\rho)(V')=z^{-|\T\setminus\T'|}\int_{(\mathfrak g)^{\mathfrak F}}\prod_{b\in\mathfrak F}\sigma(Z(b))\,dZ(b)\ \ \chi^{\rm w}\big(\hat Z(Z)\big)\prod_c\mathsf J_c(Z)\ e^{-g^{-2}\Phi_W(Z)}\,\rho\big(e^{\ii\hat Z(Z)}W\big),
\]
where $\chi^{\rm w}(\hat Z)=\prod_b\one(|\hat Z(b)|\le w)$, $\Phi_W$ is the function of Lemma~\ref{lem:exp-gauge}(c), and
\[
\mathsf J_c(Z)=\frac{\sigma(\zeta_c(Z))}{\sigma(0)\,\det\partial\ell_c/\partial Z(b_0(c))}\,,
\]
evaluated at $Z(b_0(c))=\zeta_c(Z)$. For real $Z$ this determinant is positive (see the proof), so it equals its absolute value, which is what the change of variables produces. The integrand is defined to be $0$ if $|Z(b)|>w$ for some $b\in\mathfrak F$. On the remaining set $|Z|_\infty\le w\le r_L\le c_0L^{-1}$, so $\zeta_c(Z)$, $\hat Z(Z)$, $\mathsf J_c(Z)$, $\Phi_W(Z)$ and $\sigma(Z(b))$ are defined, and the factor $\chi^{\rm w}(\hat Z(Z))$ imposes the cutoff also on the bonds $b_0(c)$. The same convention is used for all integrals over fields on $\mathfrak F$ below. Each $\mathsf J_c$ is an analytic function of $Z$ on the bonds of Lemma~\ref{lem:eliminate}(a), and $\mathsf J_c(Z)=L^{(d-1)\dim\mathfrak g}\big(1+O(L^d\mathsf a+L^{2d+1}|Z|_\infty)\big)$, with $|Z|_\infty$ the sup over $\mathfrak F$. The value on $b_0(c)$ is $\zeta_c$, with $|\zeta_c|\le CL^d|Z|_\infty$.
\end{proposition}

\begin{proof}\hasdetail{prop:exact-step}
By \eqref{eq:chi-w-support}, $\chi^{\rm w}$ vanishes outside $\mathfrak M=\mathfrak U_k(\eps_M)$, and $8\eps_ML^2=r_M$, so $\mathfrak M$ lies in the interior of $\operatorname{dom}\mathfrak a'$ (Section~\ref{sec:rg-step}). On its support $V=e^{\ii Z}W$ with $Z(b)\in\mathfrak g$, $|Z(b)|\le w$ on all bonds, and $dV(b)=\sigma(Z(b))dZ(b)$.

\emph{Submersion.} By Lemma~\ref{lem:eliminate}(a),(b), at such $V$ the derivative of $Z\mapsto(\ell_c(Z))_c$ with respect to the variables $Z(b_0(c))$ is block diagonal, with the invertible blocks $L^{1-d}R(W([c_-,b_0(c)_-]))(1+\mathfrak r_c)$; here $|\mathfrak r_c|\le\frac12$, as in the proof of Lemma~\ref{lem:eliminate}(c), since $|Z|\le w\le\rho_L$. Since $\ell_c(Z)=\ell\big(\mathfrak a'(e^{\ii Z}W)(c)\,V'(c)^{-1}\big)$, and the exponential and the logarithm are local diffeomorphisms, $\mathfrak a'$ is a submersion at $V$. So the density at $V'$ is an integral over the fibre $\{\mathfrak a'(V)=V'\}$ (Section~\ref{sec:rg-step}), which we compute with $\delta$-functions. The elimination of $Z(b_0(c))$ below is the change of variables $z\mapsto\ell_c(Z,z)$ in this integral, at fixed variables on $\mathfrak F$.

\emph{The formula.} Since $w\le r_L\le c_0L^{-1}$ for $c$ small, Lemma~\ref{lem:exp-gauge}(c) gives $\mathrm{gf}_W(V)=z^{-|\T\setminus\T'|}e^{-g^{-2}\Phi_W(Z)}$. Moreover $\mathfrak a'(V)(c)V'(c)^{-1}=\overline X(c)=e^{\ii\ell_c(Z)}$. By Lemma~\ref{lem:eliminate}(a), the $\delta$-function for $c$ involves, among the variables $Z(b_0(\cdot))$, only $Z(b_0(c))$. The function $\Phi_W$ involves none of them, since it depends only on inner bonds (Lemma~\ref{lem:exp-gauge}(c)). So the integrals over these variables factorize. For fixed variables on $\mathfrak F$, $Z(b_0(c))\mapsto\ell_c$ has at most one zero in $|Z(b_0(c))|\le w$, by Lemma~\ref{lem:eliminate}(c). It is $\zeta_c(Z)$ if $|\zeta_c(Z)|\le w$; otherwise the integrand vanishes, which is accounted for by the factor $\chi^{\rm w}(\hat Z(Z))$. Hence
\[
\int\sigma(z)\,dz\ \delta\big(e^{\ii\ell_c(Z,z)}\big)\,(\cdots)=\int\sigma(z)\frac{\delta(\ell_c(Z,z))}{\sigma(\ell_c(Z,z))}(\cdots)\,dz=\frac{\sigma(\zeta_c)}{\sigma(0)\,|\det\partial_z\ell_c|}(\cdots)\Big|_{z=\zeta_c} .
\]
\emph{Positivity of the determinant.} By Lemma~\ref{lem:eliminate}(b),
\[
\det\partial_z\ell_c=L^{(1-d)\dim\mathfrak g}\det R\big(W([c_-,b_0(c)_-])\big)\det(1+\mathfrak r_c).
\]
Here $\det R(\cdot)=1$ on $\mathfrak g$, since $G$ is connected (Section~\ref{sec:background}). In the proof of Lemma~\ref{lem:eliminate}(c), $|\mathfrak r_c|\le\frac12$ wherever $\zeta_c$ is defined, so $t\mapsto\det(1+t\mathfrak r_c)$ does not vanish on $[0,1]$, and $\det(1+\mathfrak r_c)>0$ for real $Z$. Hence $|\det\partial_z\ell_c|=\det\partial_z\ell_c$ for real $Z$, and $\mathsf J_c$ is the analytic function of the statement. The bound on $\mathsf J_c$ follows from Lemma~\ref{lem:eliminate}(b),(c), since $L^{d+1}|\zeta_c|\le CL^{2d+1}|Z|_\infty$ and $\sigma(\zeta_c)/\sigma(0)=1+O(|\zeta_c|^2)$.
\end{proof}

No invariance of $\rho$ is used: the gauge fixing has been made by the weight $\mathrm{gf}_W$, and the elimination of the constraint is a change of variables. This is Ba\l aban's representation \cite[(0.16)--(0.19)]{BalYM4I}, \cite[(2.3)--(2.7)]{BalYM4II}: the integration variables are the bonds of $\T$ other than the $b_0(c)$, and the elimination is performed on one bond per coarse bond, so that the Jacobians are local, as in \cite[(17), (20)]{BalUV3}. The representation is exact. No small-field approximation has been made beyond the insertion of $\chi^{\rm w}$, and by Lemma~\ref{lem:gf-step-invariance}(iv) the gauge fixing changes nothing on $\mathcal G'$-invariant integrands supported in $\mathfrak M$.
